\section{Conclusion}
\label{sec:conclusion}

We evaluated zero-shot GPT-4o-mini test generation on 63 HumanEval-Java functions using execution status, branch coverage, mutation score, and paired EvoSuite comparisons. GPT produced an API response for all 63 targets, but only 14 suites were executable after one permitted repair pass. Across the full corpus, branch coverage was 18.90\% and mutation score was 16.21\%; only 13 functions achieved both thresholds simultaneously. RQ1, RQ2, and RQ4 were not supported. For RQ3, no GPT--EvoSuite comparison reached significance after Holm correction on the pass-conditioned subset. Assertion failures accounted for 95.9\% of the 49 invalid suites (RQ5), identifying oracle construction as the primary obstacle in this zero-shot protocol.

Future work should pursue two directions. First, the prompt strategy should be varied systematically---comparing zero-shot, few-shot, and feedback-guided repair---while holding the SUT corpus and measurement pipeline constant, to determine whether additional context reduces the assertion-failure rate without inflating API cost. Second, the deferred student comparison requires its own protocol: student-written test suites should be collected with consent, compiled and executed in the same Maven environment, and measured without modifying their assertions. Crucially, the study must preregister how incomplete or non-compiling submissions are handled, how multiple attempts per student are counted, and whether the comparison unit is a student, a suite, or a function. Only once those choices are fixed and documented can an LLM--student statement be statistically and ethically sound. Until then, EvoSuite remains the appropriate reproducible search-based baseline, and no claim about student-written test quality should be inferred from either tool's results.
